\begin{abstract}
Taking the treatment of money and speculation in Keynes' \treatise{} and his \GT{} as a vantage point, we identify two main disagreements between these works: first, whether money is endogenous or exogenous and second, whether speculative markets feed back into real economic activity. While each book takes a distinct position on both of these questions, we argue that their recombination suitably captures key intuition prevalent in modern Post-Keynesian macroeconomics as well as centrak bank practice. We illustrate this argument by focusing on the role of speculation in shaping output and employment based on a simple model of endogenous money. We identify three distinct channels through which speculation shapes output and employment to show how the exact operations underlying speculative practices will impact on economic outcomes. We also discuss real-world implications of how our model can substantiate debates on the macroeconomic impact of speculation, potential policy-respones and impacts on prevalent financial market practices.

% TODO: Should we also add A22 for undergraduate teaching?
\noindent\textbf{JEL codes} B22, E12, E44, E51.
\end{abstract}
